% Terminal note of the homily: everything the speech itself must not carry ---
% the audience, the spoken word count and the pace it implies, the relation to
% the three readings of the expansive study, the route by which each quoted
% text reaches the page, and the exact loci --- followed by the References.
% Nothing here is read aloud.

\section*{Note on Sources and Delivery}

{\footnotesize

\textbf{Audience and occasion.} An adult parish assembly, at the Mass
\latin{Dominica decima octava post Pentecosten}, II classis, of the 1962
\work{Missale Romanum}, editio typica, pp.~410--411, nos.~1669--1678, for the
occurrence of 27~September 2026.

\textbf{Length and pace.} The spoken body runs to 1,562 words, about 12.0 to
13.0 minutes at an unhurried 120 to 130 words a minute by arithmetic alone; no
one has delivered or timed it. It was read through in full, silently, for
sense and ease of speech.

\textbf{Relation to the reviewed readings.} Of the expansive study's three
readings of the whole Mass, the homily follows the third, \textit{Thanks at Both Ends, and a Man Who Was Carried: The Grace the
Collect Confesses}: the Collect's \latin{sine te}, the Epistle's thanks for
gifts given, the man carried to Christ, Augustine's contrite heart at the
Communion, and the Postcommunion's thanks and plea. From the second,
\textit{Power on Earth to Forgive Sins: The Paralytic in His Own City and the
Altar of the Covenant}, it takes Mt 9:6 alone: the unseen forgiveness proved
by the seen healing, and Aquinas's distinction between Christ's authority and
the Apostles' ministry. Its application to Penance is the homily's own. So
are the statements that faith, sorrow, confession and the humble heart are
themselves God's gifts, which apply the Collect and 1 Cor 4:7 (which the study
quotes) and are not Augustine's gloss. The first reading,
\textit{The Peace of a City Still Being Built}, is not used. Three
disagreements the study keeps are spoken: Jerome and Ambrose against
Chrysostom on whose faith Christ saw; Chrysostom against Hilary on the crowd's
praise at Mt 9:8; and Chrysostom against Aquinas, with Ambrosiaster and
Cornelius a Lapide, on whether 1 Cor 1:8 accuses or promises. These
applications and the joins between elements, including the hearing of that
praise beside the Epistle's thanksgiving, are credited to no Father. The Introit,
Gradual, Alleluia, Offertory and Secret are left to the studies.

\textbf{Texts quoted, and their route.} Scripture is the Douay--Rheims,
Challoner's revision, at the canonical verses. The Collect and the
Postcommunion are quoted in the anonymous English printed by Eugene
Cummiskey at Philadelphia in 1861, a historical study aid and not an approved
liturgical translation of the 1962 Missal. Chrysostom on Matthew and
Augustine on Ps~95 are quoted in the English of the Nicene and Post-Nicene
Fathers, re-read in the tracked CCEL texts; Chrysostom on
First Corinthians as the study prints it from New Advent's delivery of that
series, not re-read here. Jerome, Ambrose, Hilary and Aquinas are
reported in English glosses of their Latin. Ambrose on Luke, and Hilary and
Aquinas on Matthew, were re-read in the tracked transcriptions of Migne and
the Venice optical text; Jerome on Matthew and Aquinas on First Corinthians
are reported as the study reads them.

\textbf{Exact loci.} Collect, no.~1670; Epistle, 1~Cor 1:4--8, no.~1671, with
1:10 beyond the lesson, and 4:7; Gospel, Mt 9:1--8, no.~1674, with Mk 2:4 and
Lk 5:19 for the roof; Communion, Ps 95:8--9, no.~1677; Postcommunion, no.~1678.
Chrysostom, \work{Homilies on Matthew} 29, 1 (the sick man's share in the
faith) and 2 (the healer's power, the unseen proved by the seen, and the
crowd); \work{Homilies on First Corinthians} 2, 1--5, on 1:4--5, and 7, on
1:8.
Jerome, \work{Commentarii in Matthaeum} I, on 9:1--2 (\latin{non ejus fidem qui
offerebatur, sed eorum qui offerebant}; \latin{quem sacerdotes non dignabantur
attingere}). Ambrose, \work{Expositio in Lucam} V.11 (\latin{Magnus Dominus qui
aliorum merito ignoscit aliis}; \latin{adhibe Ecclesiam quae pro te
precetur}). Aquinas, \work{Super Matthaeum} c.~9 (\latin{per viam
administrationis, non auctoritatis}; \latin{quia portabatur, praecepit ut
portaret}; \latin{quia ire non poterat, dixit, Et ambula}), and \work{Super I ad Corinthios} c.~1, lect.~1, on 1:8. Hilary,
\work{Commentarius in Matthaeum} VIII.8 (\latin{potestas hominibus ac via data
sit per verbum ejus}). Augustine, \work{Enarrationes in Psalmos} 95, 9.

\par}

\Needspace{14\baselineskip}
\section*{References}

{\footnotesize
\begin{itemize}\setlength{\itemsep}{0.15em}\setlength{\parskip}{0pt}
\item \work{Missale Romanum ex decreto Sacrosancti Concilii Tridentini restitutum, Summorum Pontificum cura recognitum}, editio typica (Typis Polyglottis Vaticanis, 1962), CMAA facsimile, pp.~410--411, nos.~1669--1678.
\item \work{The Roman Missal, Translated into the English Language for the Use of the Laity} (Philadelphia: Eugene Cummiskey, 1861), pp.~444--445 (XVIII Sunday after Pentecost, Collect and Postcommunion).
\item The Holy Bible, Douay--Rheims version, Challoner's revision: Ps 95:8; Mt 9:1--8; Mk 2:4; Lk 5:19; 1 Cor 1:4--10, 4:7.
\item St John Chrysostom, \work{Homilies on Matthew} 29, trans. G. Prevost, rev. M.~B. Riddle, Nicene and Post-Nicene Fathers, 1st ser., vol.~10; \work{Homilies on First Corinthians} 2, 1st ser., vol.~12 (New Advent).
\item St Jerome, \work{Commentarii in Matthaeum} I, on 9:1--2 (Migne, PL 26, cols.~54--55).
\item St Ambrose, \work{Expositio Evangelii secundum Lucam} V.11 (Migne, PL 15).
\item St Hilary of Poitiers, \work{Commentarius in Matthaeum} VIII.8 (Migne, PL 9, cols.~961--962).
\item St Thomas Aquinas, \work{Super Evangelium S.~Matthaei lectura} c.~9 (Venice, 1745); \work{Super I ad Corinthios} c.~1, lect.~1 (Corpus Thomisticum).
\item St Augustine, \work{Enarrationes in Psalmos} 95, 9, Nicene and Post-Nicene Fathers, 1st ser., vol.~8.
\end{itemize}
\par}
